\section{Finite affine potentials and prescribed cell masses}
\label{sec:cells}

We return to an arbitrary compact convex body $K\subset\R^d$ and its
uniform probability measure $\rho=\mathbf 1_K\,dx/|K|$. To compare targets
by moving their sites, we need intercepts that keep prescribed cell masses
fixed and depend differentiably on the sites. We first compute the derivatives
of cell masses and moments, then use them to construct these intercepts.

Assume that the sites $y_1,\ldots,y_N$ are distinct. For intercepts
$c=(c_1,\ldots,c_N)$, set
\begin{align*}
 P_i&=K\cap\bigcap_{j\ne i}
 \{x:x\cdot(y_i-y_j)+c_i-c_j\ge0\},\\
 M_i&=\rho(P_i),
 &B_i&=\int_{P_i}x\,d\rho(x).
\end{align*}
These closed convex cells partition $K$ up to the null tie set. Denote the
equality plane and the common face by
\[
 H_{ij}=\{x:x\cdot(y_i-y_j)+c_i-c_j=0\},
 \qquad F_{ij}=P_i\cap P_j.
\]
All integrals over $F_{ij}$ use the $(d-1)$-dimensional measure $dS$ of
$H_{ij}$, so an empty or lower-dimensional face contributes zero. Define
the symmetric coefficients
\begin{equation}
\label{cells:face-coefficients}
 g_{ij}=\frac{1}{|K|\,|y_i-y_j|}\int_{F_{ij}}1\,dS,
 \qquad
 b_{ij}=\frac{1}{|K|\,|y_i-y_j|}\int_{F_{ij}}x\,dS.
\end{equation}

These cells are often called Laguerre cells. Their fixed-site mass
derivative has the standard face-integral and weighted graph-Laplacian
form \cite[Theorem~1.3 and Section~5.3]{KitagawaMerigotThibert2019}.
We include the calculation for the joint site and moment variations
needed here.

\begin{proposition}[Cell differentials and connectivity]
\label{cells:calculus}
On the open set of configurations with distinct sites and $M_i>0$ for
every $i$, the maps $M_i$ and $B_i$ are jointly $C^1$ in the intercepts and
sites. For an intercept direction $r=(r_i)$ and a site direction
$v=(v_i)$,
\begin{align}
 D M_i[r,v]
 &=\sum_{j\ne i}\bigl[g_{ij}(r_i-r_j)
          +b_{ij}\cdot(v_i-v_j)\bigr],
 \label{cells:mass-differential}\\
 D B_i[r,v]
 &=\sum_{j\ne i}b_{ij}(r_i-r_j)
   +\frac1{|K|}\sum_{j\ne i}\frac1{|y_i-y_j|}
        \int_{F_{ij}}x\bigl(x\cdot(v_i-v_j)\bigr)\,dS(x).
 \label{cells:moment-differential}
\end{align}
The graph on $\{1,\ldots,N\}$ with edges $g_{ij}>0$ is connected.
Consequently, the symmetric operator
\begin{equation}
\label{cells:laplacian}
 (Gq)_i=D M_i[q,0]=\sum_{j\ne i}g_{ij}(q_i-q_j)
\end{equation}
has kernel consisting of the constant vectors and range
$\{z\in\R^N:\sum_i z_i=0\}$.
\end{proposition}

\begin{proof}
We first record the geometric facts that permit differentiation for an
arbitrary convex body. The boundary of a full-dimensional compact convex
set has zero volume: contractions toward an interior point lie in the
interior and their volumes tend to the full volume. The same argument
applies within an affine subspace. Moreover, if a hyperplane $H$ meets
$\operatorname{int}K$, then
\[
 \operatorname{relint}(K\cap H)=\operatorname{int}K\cap H.
\]
Indeed, choose $p\in\operatorname{int}K\cap H$. A relative interior point
$z\ne p$ of the section lies strictly between $p$ and a point of the
section beyond $z$. A strict convex combination with the ambient interior
point $p$ lies in $\operatorname{int}K$. The reverse inclusion is immediate.
It follows that $\partial K\cap H$ has zero $(d-1)$-dimensional measure.

Since $M_i>0$, cell $i$ contains an interior point of $K$ at which its
affine function is strictly maximal: remove the boundary of $K$ and the
finite union of tie planes from the positive-volume cell. Strictly active
points for $i$ and $j$ give opposite signs of their affine difference, so
$H_{ij}$ meets $\operatorname{int}K$. Also, $H_{ij}$ and $H_{ik}$ cannot
coincide for three distinct indices. Coincidence would make their affine
differences proportional. The three affine functions would then lie on
one affine line, with one a strict convex combination of the other two;
that function could never be strictly maximal. Hence the intersections
$H_{ij}\cap H_{ik}$ are empty or have dimension at most $d-2$.

We next prove the differential formulas. For a continuous function
$f:\R^d\to\R$, write
\[
 I_i^f=\frac1{|K|}\int_K f(x)
        \prod_{j\ne i}\mathbf1_{\{\ell_{ij}(x)\ge0\}}\,dx,
 \qquad
 \ell_{ij}(x)=x\cdot(y_i-y_j)+c_i-c_j.
\]
Only $f=1$ and $f(x)=x_k$ will be needed. Perturb $(c,y)$ by $h(r,v)$,
with $h\downarrow0$, and telescope the difference of the products of
indicators one factor at a time. On bounded $K$, the term in which the
$j$-th factor changes is supported in
$|\ell_{ij}(x)|\le Ch$, with $C$ fixed for the configuration and direction.
Put $n_{ij}=(y_i-y_j)/|y_i-y_j|$ and use coordinates
\[
 x=z+htn_{ij},\qquad z\in H_{ij}.
\]
Both $z$ and $t$ can be restricted to fixed bounded sets. The factor $h$
in the volume element cancels the difference quotient. Outside
\[
 (\partial K\cap H_{ij})
 \ \cup\!\bigcup_{k\ne i,j}(H_{ik}\cap H_{ij}),
\]
a surface-null set by the preceding geometric observations, all remaining
indicators converge to their values at $z$. The changed factor converges
for almost every $(z,t)$ to
\[
 \mathbf1_{\{|y_i-y_j|t+r_i-r_j+z\cdot(v_i-v_j)\ge0\}}
       -\mathbf1_{\{t\ge0\}}.
\]
Its integral in $t$ is
$(r_i-r_j+z\cdot(v_i-v_j))/|y_i-y_j|$. Dominated convergence therefore gives
\begin{equation}
\label{cells:general-differential}
 D I_i^f[r,v]
 =\frac1{|K|}\sum_{j\ne i}\frac1{|y_i-y_j|}
    \int_{F_{ij}} f(z)
       \bigl(r_i-r_j+z\cdot(v_i-v_j)\bigr)\,dS(z).
\end{equation}
Applying the same argument in the opposite direction supplies the
two-sided derivative. Taking $f=1,x_1,\ldots,x_d$ gives
\eqref{cells:mass-differential} and \eqref{cells:moment-differential}.

For completeness, these derivatives are continuous under joint changes
of sites and intercepts. Identify a nearby equality plane with the fixed
$H_{ij}$ as an affine graph in the normal direction $n_{ij}$. The graph
maps and their surface Jacobians converge uniformly on a common bounded
integration region. The face-defining indicators converge outside the
same surface-null sets used above. The remaining integrands and the
positive denominators $|y_i-y_j|$ converge as well. Dominated convergence
proves continuity of each coefficient of the differential, including
those containing $z_kz_l$. Thus all coordinate partial derivatives exist
and are continuous, which yields joint $C^1$ regularity. The admissible
parameter set is open: the masses themselves are continuous by dominated
convergence away from the original tie hyperplanes.

To prove connectivity, fix strictly active interior points in any two
cells. They may be chosen so that the first point $p$ avoids all equality
planes and the second point $q$ avoids all equality planes and the affine
spans of $p$ with every triple equality set. These are finitely many
proper affine subsets, so they cannot cover a strictly active open
region. The segment $[p,q]$ lies in $\operatorname{int}K$ and meets no
triple equality set. At any change of active index along it, precisely
two affine functions are maximal. The other inequalities are strict,
so a small relatively open disk in their equality plane belongs to the
common face. That face has positive area and gives an edge of the graph.
The successive active indices therefore give a graph path between the
chosen cells.

Finally, symmetry of the face coefficients gives
\[
 q\cdot Gq=\sum_{i<j}g_{ij}(q_i-q_j)^2.
\]
Connectivity identifies the kernel as the constant vectors. Symmetry
then identifies the range as their orthogonal complement, as claimed.
\end{proof}

\begin{proposition}[Prescribed cell masses]
\label{cells:weights}
Fix positive weights $m_1,\ldots,m_N$ with $\sum_i m_i=1$. For every
configuration of distinct sites, there is a unique vector of intercepts
$c$ satisfying
\begin{equation}
\label{cells:prescribed-masses}
 M_i=m_i\quad(1\le i\le N),\qquad \sum_i m_i c_i=0.
\end{equation}
With the weights fixed, this normalized vector is a $C^1$ function of
the sites on the open set of distinct-site configurations.
\end{proposition}

\begin{proof}
The case $N=1$ has $c_1=0$, so assume $N\ge2$. On the normalization space
$H_m=\{c\in\R^N:\sum_i m_i c_i=0\}$, minimize the convex function
\[
 \mathcal F_y(c)=\int_K\max_i(x\cdot y_i+c_i)\,d\rho(x)
                  -\sum_i m_i c_i.
\]
Choose $R$ with $|x|\le R$ on $K$, and set $S=\max_i|y_i|$. On $H_m$,
\[
 \mathcal F_y(c)\ge\max_i c_i-RS.
\]
Positivity of all the weights implies that $\max_i c_i$ tends to infinity
as $|c|\to\infty$ within $H_m$: an upper bound on all coordinates and the
weighted mean-zero identity also bound every coordinate below. Hence
$\mathcal F_y$ is coercive on $H_m$ and attains its minimum.

Distinct sites give an almost-everywhere unique maximizing index for
every $c$. Differentiation of the maximum under the integral, justified
by bounded difference quotients, shows
\[
 \nabla_c\mathcal F_y=(M_i-m_i)_{i=1}^N.
\]
At a minimum on $H_m$, this gradient lies in $\operatorname{span}\{m\}$.
Its coordinate sum is zero because the masses and weights each sum to
one. Since $\sum_i m_i=1$, the gradient must vanish. Thus a normalized
solution of \eqref{cells:prescribed-masses} exists. Conversely, every such
solution is a minimizer by convexity.

At any solution all cells have positive mass, so \Cref{cells:calculus}
applies. The intercept derivative of the mass residual is the restriction
\[
 G:H_m\longrightarrow H_0,
 \qquad H_0=\{z\in\R^N:\textstyle\sum_i z_i=0\}.
\]
It is an isomorphism: both spaces have dimension $N-1$, and the only
constant vector in $H_m$ is zero. The finite-dimensional implicit
function theorem therefore gives a locally unique $C^1$ solution as the
sites vary. In particular, each solution is locally isolated for its
fixed sites. If two normalized solutions existed, their segment would
consist of minimizers of the convex objective, contradicting this local
isolation. Global uniqueness follows, and the local $C^1$ solution
branches agree wherever they overlap. This proves the claimed $C^1$
dependence on the sites. Only the invertibility of the restricted
Laplacian is used here; no quantitative bound on its inverse is needed.
\end{proof}
